\documentclass[conference]{IEEEtran}
\usepackage{cite}
\usepackage{amsmath,amssymb,amsfonts}
\usepackage{booktabs}
\usepackage{array}
\usepackage{multirow}
\usepackage{url}
\usepackage{xcolor}
\newcommand{\indicator}{\mathbb{I}}
\begin{document}

\title{Will Humans Look Physically Different in the Future?\\
A Pathway-Attribution Proposal for Population-Level Physical Phenotypes}

\author{\IEEEauthorblockN{Anonymous Research Plan Author}\\
\IEEEauthorblockA{Source-bound Survey--Idea--ExperimentDesign research proposal\\
Design-only scientific study; no observed results are reported}}

\maketitle

\begin{abstract}
Claims about how future humans will look are often presented as if visible
change had one cause: evolution. That framing is scientifically inadequate.
A difference between population cohorts can reflect transgenerational
allele-frequency change, developmental plasticity, technology or treatment,
demographic composition, or a change in measurement practice. These pathways
have different time scales, causal meanings, data requirements, and ethical
risks. This paper converts the broad question into a falsifiable research
proposal: can a preregistered pathway-aware model provide a better calibrated
and more transportable account of selected population-level physical
phenotypes than an age, birth-cohort, site, and measurement-process baseline?
The proposed Phenotype Pathway Attribution Atlas is a noninterventional
longitudinal study using only eligible, consented records. It analyzes
pre-specified phenotype families separately, including congenital third-molar
agenesis as a bounded dental case study, de-identified craniofacial scalar
measurements, and harmonized anthropometry. The protocol prevents a common
error in the wisdom-tooth narrative by separating congenital agenesis from
eruption, impaction, extraction, and treatment. Optional inherited-feature
summaries are subject to separate consent and stringent population-structure
sensitivity analyses; they are never used to rank individuals or human
groups. Temporal and external-site holdouts, calibration tests,
measurement-era ablations, negative controls, and governance release gates
define what would count as support. The proposal does not forecast a universal
future appearance, does not claim that any trend has been observed, and does
not make clinical or reproductive recommendations. Its contribution is an
auditable way to replace visually compelling but unsupported speculation with
conditional, population-level scientific inference.
\end{abstract}

\begin{IEEEkeywords}
human phenotype, developmental plasticity, population genetics, third-molar
agenesis, cohort effects, prediction calibration, demographic composition,
research governance
\end{IEEEkeywords}

\section{Introduction}

It is reasonable to expect that human populations in the future will not be
identical to those observed today. It is not reasonable to convert that
general possibility into a drawing of a universal future face, body, or
dentition. Physical phenotype is the result of development within an
environment, inherited variation, disease and treatment history, material and
social conditions, and the way bodies are sampled and measured. A cohort
difference can therefore be real while its popular explanation is wrong. This
distinction is especially important when a claim implicitly treats geographic,
ancestry, or demographic categories as fixed biological types.

The term evolution is frequently used too broadly. In a population genetic
sense, evolution requires a change in inherited variant frequencies across
generations. In contrast, developmental plasticity denotes phenotypic change
during ontogeny in response to biological and environmental conditions.
Technology may alter phenotype through dentistry, orthodontics, medicine,
nutrition, assistive devices, or surgery without making that phenotypic change
inherited. Migration and changing cohort composition can alter a population
average with no within-lineage biological transformation. Finally, a new
scanner, clinical referral rule, or measurement protocol can make a trend
appear even when the underlying distribution is unchanged.

The prompt's reference to wisdom teeth is a useful illustration of the
problem. Third-molar agenesis is a congenital developmental outcome, while
wisdom-tooth eruption, impaction, extraction, and orthodontic management are
different outcomes. Carter \emph{et al.} synthesized 92 studies comprising
63,314 subjects and estimated a heterogeneous worldwide agenesis rate of
22.63 percent, with estimates ranging from 5.32 to 56.0 percent
\cite{carter2015}. This is evidence that agenesis is a measurable and
heterogeneous phenotype. It is not evidence that every population is
currently evolving toward tooth loss, that a smaller jaw causes congenital
agenesis, or that the frequency must rise in the future.

The same caution applies to genetic forecasts. Modern genomic data make it
possible to examine associations with complex traits, but association does
not remove population structure, environmental correlation, or
transferability problems. Berg \emph{et al.} reported that previously
prominent polygenic-adaptation signals for height were strongly attenuated or
absent when reexamined with independent UK Biobank estimates, and documented
confounding by population stratification \cite{berg2019}. Accordingly,
polygenic-score contrasts are not a safe substitute for causal or predictive
evidence about visible human differences.

This paper presents a design-only proposal rather than a claim of completed
research. It asks a deliberately bounded question: \emph{after outcome
definitions and analysis decisions have been locked, does a longitudinal
model that represents demographic composition, developmental exposure,
technology or treatment, measurement process, and optional inherited features
separately improve calibration and transportability over a cohort-only
baseline?} This is a question about population-level research inference. It
does not predict an individual appearance and it does not authorize a
conclusion about a group merely because the group appears in a dataset.

\section{Background, Survey, and Research Gap}

\subsection{Pathways that must not be merged}

Human developmental systems can respond to conditions during ontogeny.
Ponzi \emph{et al.} review hormonal mediation of developmental plasticity and
emphasize that changes in phenotype can arise during development through
endocrine systems responsive to environmental conditions \cite{ponzi2020}.
That evidence supports a modeling distinction: a population trend in an
anthropometric or craniofacial measure may be associated with developmental
exposures without representing an inherited change in allele frequency. The
source does not specify a universal direction for any future human feature.

At longer time scales, human cultural practices can affect environments in
which genetic evolution occurs. Laland \emph{et al.} describe gene--culture
coevolution as an important framework for connecting cultural activities,
selection environments, and human genetic change \cite{laland2010}. This
permits technology, diet, and other cultural variables to be considered as
candidate environmental pathways. It does not justify treating a present-day
technology association as proof of selection or as a forecast of an inherited
trait.

Population composition is a separate pathway. If age distribution, parental
birthplace distribution, migration pattern, recruitment source, or survival
changes over time, a cohort mean can change even when no individual process
has changed. This requires explicit covariates, site effects, and temporal
validation. The 1000 Genomes Project reference is a reminder that genotype
data require careful quality control and documented structure rather than
superficial biological labels \cite{auton2015}.

The final pathway is measurement and intervention. Orthodontic treatment can
affect dental alignment; extraction removes a tooth after development; medical
care and nutrition can affect developmental conditions; and a new radiographic
or imaging protocol can alter ascertainment. These variables may be important
predictors, but a predictive association is not a causal effect estimate.
Without a separately valid causal design, this study will describe
associations and forecast performance only.

\begin{table*}[t]
\caption{Pathways that can produce apparent physical differences and the
inference each pathway permits}
\label{tab:pathways}
\centering
\footnotesize
\begin{tabular}{p{0.19\textwidth}p{0.27\textwidth}p{0.22\textwidth}p{0.22\textwidth}}
\toprule
Pathway & Example observable inputs & Permitted inference in this study & Prohibited shortcut \\
\midrule
Inherited population change & Separately consented, quality-controlled genotype summaries; family and structure information & A candidate predictive block may be evaluated after ancestry-aware sensitivity analyses. & Treating a polygenic-score contrast as a ranking of groups or a deterministic appearance forecast. \\
Developmental plasticity & Time-aligned developmental exposures and health or nutrition indicators where consented & A developmental context can be associated with a phenotype and may improve prediction. & Calling an ontogenetic response inherited evolution. \\
Technology and treatment & Dental treatment and extraction history, measurement-appropriate care access proxies & Treatment history can separate congenital from post-developmental outcomes. & Calling an observational treatment coefficient causal. \\
Demographic composition & Eligibility, recruitment, site, calendar period, and composition descriptors & A cohort average can be decomposed and audited for compositional change. & Treating a changing sample as a changing biological type. \\
Measurement process & Instrument, protocol, assessor, referral pathway, and measurement era & Method changes can be modeled or can invalidate an endpoint. & Calling a protocol artifact a biological trend. \\
\bottomrule
\end{tabular}
\end{table*}

\subsection{Bounded dental endpoint and genetic caution}

The third molars are the most frequently affected teeth in congenital
agenesis, but this fact should increase, not reduce, measurement discipline.
The meta-analysis by Carter \emph{et al.} shows considerable geographic and
demographic heterogeneity \cite{carter2015}. A record of no visible third
molar is not automatically a record of agenesis. The tooth may not yet have
erupted, may be impacted, may have been extracted, or may have an
unrecoverable history. These alternatives carry different biological and
clinical meanings.

The proposed dental endpoint therefore requires an age-adequate developmental
assessment and a provenance field that explicitly records congenital absence,
known extraction, impaction, eruption status, and treatment history. If a
source cannot support that separation, it cannot contribute to the
confirmatory dental analysis. This conservative choice can reduce sample size,
but it prevents a visually memorable narrative from outweighing endpoint
validity.

The proposed study does not deny genetic contribution to physical phenotype.
It refuses to assume that such contribution can be read directly from
population labels or from a score optimized in another sample. Berg
\emph{et al.} provide a direct warning: previous height-adaptation results
were susceptible to residual structure and should be interpreted cautiously
\cite{berg2019}. A genotype summary may therefore enter the design only after
separate consent, documented quality control, relatedness-aware procedures,
and analysis that tests whether the finding survives reasonable choices about
structure. If it does not, the block is removed. A simpler model with
transparent limits is preferable to an impressive model that encodes
unresolved confounding.

\subsection{Frozen research gaps}

The Survey freezes three linked gaps. \textbf{GAP-FP-01} is pathway
conflation: future-phenotype claims often merge genetic change, plasticity,
technology, composition, and measurement. \textbf{GAP-FP-02} is outcome
conflation in dental narratives: congenital agenesis is commonly discussed
beside extraction and jaw-related outcomes without a common endpoint
definition. \textbf{GAP-FP-03} is validation weakness: cohort-based forecasts
often omit temporal transport, external-site transport, calibration, and
sensitivity to population structure or measurement change. The project is
designed to address these gaps, not to answer every question about human
evolution.

\section{Research Questions and Planned Contributions}

The primary research question is: \emph{Does a preregistered pathway-aware
longitudinal model improve out-of-sample calibration and proper scoring for
one or more selected population-level physical phenotypes beyond an age,
birth-cohort, site, and measurement-process baseline?}

The supporting questions are whether developmental variables improve
prediction after the baseline is locked, whether treatment records resolve
endpoint misclassification, whether gains persist across later periods and
independent sites, and whether any inherited-feature block remains stable
after structure and relatedness sensitivity analyses. The planned
contributions are a pathway taxonomy, a strict dental outcome definition, a
common baseline, locked transport tests, and a communication boundary against
individual appearance prediction and population ranking.

\section{Idea Source Checkpoints and Direction Selection Audit}

The Idea stage compared three directions. A Secular Phenotype Trend Catalogue
would inventory cohort differences in teeth, faces, and anthropometry. It is
useful descriptively but cannot attribute a difference to a distinct pathway.
It was therefore not selected. The selected \emph{Phenotype Pathway
Attribution Atlas} uses longitudinal records, separate predictor blocks, and
falsification gates. It addresses all three frozen gaps. A universal future-
human visual generator was rejected because it has no observable universal
target and would produce deterministic, stereotype-prone claims.

No Monte Carlo tree search, autonomous evolutionary optimization, or
simulated multi-agent debate was executed. The selection record is
transparent rather than fictional. The selected idea fails if pathway-aware
models do not improve locked holdout performance, if improvements disappear
after endpoint harmonization or sensitivity analysis, or if governance gates
cannot be met.

\section{Problem Definition, Assumptions, and Hypotheses}

\subsection{Target population and endpoints}

Let $Y_{ijk}$ denote phenotype $k$ for person $i$ at visit or age-interval
$j$. Each endpoint has its own eligibility rule and is analyzed separately.
The first confirmatory endpoint is congenital third-molar agenesis, denoted
by $A_i$. It is defined only when an age-adequate source record supports a
developmental determination:
\begin{equation}
A_i=\indicator\{\text{congenital third-molar absence is documented}\}.
\label{eq:agenesis}
\end{equation}
The indicator in \eqref{eq:agenesis} is not assigned from an absent tooth in
a later clinical record. Known extraction, impaction, eruption status, and
treatment are separate variables. This does not prevent descriptive analyses
of those outcomes; it prevents them from entering the same outcome label.

The second family comprises de-identified craniofacial scalar measures, such
as pre-registered linear distances or ratios, available only from approved
existing data. The third family includes standardized anthropometric measures
where an instrument and protocol can be harmonized or modeled. Raw facial
images, free-text appearance labels, and subjective attractiveness ratings are
excluded. The design has no requirement that all phenotype families be
available at every site; the prespecified cohort eligibility table will
identify which sites support which endpoint.

\subsection{Model classes}

The baseline model represents age or developmental stage, birth cohort, site,
and measurement process. It intentionally includes variables that can generate
a spurious-looking trend:
\begin{equation}
\eta^{(0)}_{ijk}=\alpha_k+g_k(t_{ij})+\gamma_{k,c(j)}
+u_{k,s(i)}+\delta_{k,m(ij)}.
\label{eq:baseline}
\end{equation}
In \eqref{eq:baseline}, $\eta^{(0)}_{ijk}$ is the baseline linear predictor
for phenotype $k$, $g_k(\cdot)$ is a preregistered flexible function of age
or developmental stage $t_{ij}$, $\gamma_{k,c(j)}$ is the birth-cohort
effect, $u_{k,s(i)}$ is the site effect, and $\delta_{k,m(ij)}$ represents
the measurement process. The chosen link function is binary or continuous as
appropriate to the endpoint.

The pathway-aware model augments, but does not replace, that baseline:
\begin{align}
\eta^{(1)}_{ijk}={}&\eta^{(0)}_{ijk}
+\boldsymbol{x}_{ij}^{\mathsf{T}}\boldsymbol{\beta}_{k}
+\boldsymbol{z}_{ij}^{\mathsf{T}}\boldsymbol{\theta}_{k}\nonumber\\
&+\boldsymbol{q}_{ij}^{\mathsf{T}}\boldsymbol{\phi}_{k}
+\boldsymbol{r}_{i}^{\mathsf{T}}\boldsymbol{\psi}_{k}.
\label{eq:pathway}
\end{align}
Here, $\boldsymbol{x}_{ij}$ is the demographic-composition block,
$\boldsymbol{z}_{ij}$ the developmentally time-aligned exposure block,
$\boldsymbol{q}_{ij}$ the technology or treatment block, and
$\boldsymbol{r}_{i}$ an optional inherited-feature summary. The associated
coefficient vectors are $\boldsymbol{\beta}_{k}$,
$\boldsymbol{\theta}_{k}$, $\boldsymbol{\phi}_{k}$, and
$\boldsymbol{\psi}_{k}$. Equation \eqref{eq:pathway} is a predictive,
multivariable representation. It is not a causal structural equation; in
particular, $\boldsymbol{\phi}_{k}$ is not a treatment-effect estimator.

The quantity evaluated at a future or held-out observation is a conditional
predictive distribution:
\begin{equation}
\widehat{p}_{ijk}=\Pr(Y_{ijk}\mid\mathcal{H}_{ij},
\mathcal{M}^{(1)}_k),
\label{eq:prediction}
\end{equation}
where $\mathcal{H}_{ij}$ is the record history allowed by the locked
prediction horizon and $\mathcal{M}^{(1)}_k$ is the pathway-aware model for
endpoint $k$. All imputation, transformations, feature selection, and
hyperparameter choices are trained without opening the final evaluation data.

\subsection{Assumptions and hypotheses}

The study assumes that at least one eligible source can document phenotype
measurement provenance, calendar period, site, and the core baseline
variables. It does not assume that all sites measure the same trait with the
same instrument. Instead, harmonization is an eligibility and modeling
requirement. Missing exposure data are not assumed to be random; the protocol
requires a missingness description, complete-case comparison, and a
preregistered imputation sensitivity analysis.

\textbf{H1} states that, for at least one pre-specified endpoint, the
pathway-aware model will improve a proper score and probability calibration on
both a temporal and an external-site holdout. \textbf{H2} states that
separating congenital agenesis from treatment and post-developmental dental
outcomes will materially alter descriptive interpretation compared with an
undifferentiated missing-wisdom-tooth label. Neither hypothesis predicts the
sign of a cohort trend. Both fail if the registered gates are not met.

\section{Study Design and Methods}

\subsection{Data eligibility, consent, and data minimization}

The protocol is noninterventional. It may use only existing longitudinal
cohorts or prospectively consented collections with a clear ethics approval,
lawful data-use basis, and a specific prohibition on individual appearance
prediction. The final source list, participant count, age range, follow-up
duration, and phenotype coverage are human-review inputs. They are not
invented in this proposal.

The data custodian must supply a data dictionary that includes endpoint
provenance, measurement protocol, calendar period, site, referral or sampling
pathway, and missingness indicators. Direct identifiers, raw facial images,
and free-text descriptions of appearance are not analytic deliverables.
Craniofacial data, when eligible, are reduced to pre-specified de-identified
scalar measurements in an approved secure environment. Genotype data, if
considered, require separate consent and are never exported as a score tied to
an identifiable person.

\subsection{Phenotype adjudication and measurement harmonization}

An endpoint-adjudication manual is completed before data analysis. For the
dental case study, a clinician with relevant expertise independently reviews
the coding rules and a random, preregistered sample of eligible source records
within the secure setting. Disagreements are recorded, not silently resolved.
The protocol does not ask for radiation or imaging beyond the recorded care
pathway.

For anthropometry, the manual specifies instrument make, calibration
procedure, posture, time of day where available, units, and allowable
transformation. For craniofacial scalars, it specifies landmarks, measurement
method, version, and a tolerance for inter-rater or algorithm-versus-human
agreement. If an instrument change produces an unresolvable discontinuity,
the affected records are separated, modeled as a different measurement era,
or excluded from confirmatory analysis. A larger dataset is not worth an
invalid target.

\begin{table*}[t]
\caption{Pre-analysis register for phenotype families and pathway blocks}
\label{tab:register}
\centering
\footnotesize
\begin{tabular}{p{0.20\textwidth}p{0.27\textwidth}p{0.23\textwidth}p{0.20\textwidth}}
\toprule
Element & Required operational record & Primary purpose & Exclusion or downgrade trigger \\
\midrule
Third-molar agenesis & Age-adequate developmental evidence; extraction, eruption, impaction, and treatment fields & Bounded congenital dental outcome & Absence cannot be separated from post-developmental tooth status. \\
Craniofacial scalar family & Approved de-identified measurement definitions; landmark and protocol version & Quantitative morphology without releasing faces & Raw images are needed, landmark method changes without bridge data, or agreement falls below the registered tolerance. \\
Anthropometric family & Instrument, units, protocol, and measurement era & Test a standard bodily measure with explicit process adjustment & Unit or protocol cannot be harmonized or modeled. \\
Developmental block & Time-aligned, consented exposure data with provenance and missingness fields & Separate ontogenetic context from cohort labels & Timing or consent is insufficient. \\
Technology and treatment block & Dental and health treatment history or documented access proxy & Prevent treatment from masquerading as congenital biology & Variable cannot distinguish pre- and post-outcome timing. \\
Optional inherited block & Separate consent, genotype QC, relatedness and structure records & Sensitivity-tested incremental prediction only & Structure, transferability, or subgroup calibration analysis is inadequate. \\
\bottomrule
\end{tabular}
\end{table*}

\subsection{Model development and leakage prevention}

For each endpoint, the baseline \eqref{eq:baseline} is first fit using
training records. The pathway-aware model \eqref{eq:pathway} is then fit with
exactly the pre-specified blocks available before the prediction horizon. Any
preprocessing operation that learns from a distribution, including
normalization, imputation, missingness-indicator construction, and feature
selection, is fit within the training fold and applied to held-out data. Site
and family identifiers are used to prevent the same participant or close
family group, where known, from appearing on both sides of an evaluation
split. This avoids artificial performance created by record overlap.

The model class is chosen according to endpoint type. A binary dental
endpoint can use a hierarchical logistic model or an equally transparent
pre-registered alternative. A continuous scalar endpoint can use a
hierarchical continuous-outcome model with residual and interval calibration.
Complex machine-learning alternatives are permitted only when they use the
same target, time horizon, predictor availability, and splits as the
baseline. They cannot obtain an apparent gain by reading variables recorded
after the outcome or by exploiting a source-specific identifier.

The optional genetic block is evaluated only after a non-genetic
pathway-aware model has been analyzed. Its contribution is reported as an
incremental comparison and must be accompanied by sensitivity to ancestry
structure, relatedness, covariate specification, and site composition. A
failure of this block does not invalidate the broader study; it demonstrates
that the block should not be used. This conforms to the reason for including
the block as optional in the first place: a potentially informative source of
variation is not automatically a transportable or interpretable predictor.

\subsection{Validation, calibration, and uncertainty}

The primary evaluation is a locked temporal holdout, such as a later
measurement era or later birth cohort selected before model fitting. The
secondary evaluation is an external-site holdout. Both are essential because
a model that describes one historical source may fail when an instrument,
care setting, population composition, or exposure distribution changes.
Separate within-site random splitting is insufficient for the scientific
question because it can preserve the same measurement quirks on both sides of
the split.

For a binary endpoint with outcome $A_i$ and predicted probability
$\widehat{p}_i$, one registered proper score is the Brier score:
\begin{equation}
\mathrm{BS}=\frac{1}{n}\sum_{i=1}^{n}(A_i-\widehat{p}_i)^2.
\label{eq:brier}
\end{equation}
In \eqref{eq:brier}, $n$ is the size of the locked evaluation set. Lower
scores are better, but a difference is interpretable only with an uncertainty
interval and comparison against the baseline. Calibration intercept and slope
are reported with the reliability display. For continuous endpoints, the
protocol uses an outcome-appropriate proper score, residual calibration, and
predictive-interval coverage. A ranking metric alone is not accepted because
good ranking can coexist with systematically overconfident predictions.

TRIPOD motivates complete disclosure of model development and validation
\cite{collins2015}; PROBAST supplies a structured audit of participants,
predictors, outcomes, and analysis \cite{wolff2019}. In this protocol,
these standards are not treated as badges. Their logic is embedded in the
locked target, the data dictionary, the split plan, and the release gates.

\subsection{Sample support, missingness, and sensitivity planning}

No numerical sample size is fabricated before the eligible cohorts are
known. Instead, the statistical review pre-registers an outcome-specific
simulation plan once the prevalence, number of sites, correlation structure,
and allowed degrees of freedom are documented. The plan estimates the
precision of calibration and proper-score differences under the actual
candidate data geometry. If an endpoint or a prespecified stratum lacks
support for a meaningful calibration assessment, the result is labeled
insufficiently supported rather than reported with a falsely precise
estimate.

Missing data are potentially informative. An absent treatment history may
mean that no treatment occurred, that care was delivered outside the record,
or that the field was not collected during that era. The analysis records
which of these mechanisms is plausible for every variable. It compares a
complete-case analysis with a preregistered imputation approach and with
missingness indicators where justified. A missingness procedure is not
allowed to use held-out outcomes. Material disagreement among the analyses is
reported as a limitation, not hidden by selecting the most favorable result.

The robustness package includes leave-one-site-out evaluation, leave-one-
measurement-era-out evaluation, removal of each pathway block, and negative
controls when a defensible control can be defined. For example, a purported
developmental-window exposure should not show the same association with a
temporally impossible proxy. A negative control does not identify the causal
effect of an exposure; it can reveal that an association is compatible with
measurement or composition artifact.

\subsection{Reproducibility and preregistration}

Before opening the final temporal or site holdout, the study archives a
versioned analysis plan, data dictionary, endpoint adjudication manual,
feature availability table, split definitions, performance metrics, and
interpretation rules. Code is executed only in the approved analysis
environment. It produces a lineage record for input versions, transformations,
model parameters, and tables. No raw facial image, person-level genotype, or
identifiable record is included in a public replication package. A synthetic
or metadata-only demonstration may be released only after governance review
confirms that it cannot be mistaken for observed study results.

\begin{table}[t]
\caption{Release gates for any conditional forecast statement}
\label{tab:gates}
\centering
\footnotesize
\begin{tabular}{p{0.29\columnwidth}p{0.61\columnwidth}}
\toprule
Gate & Requirement \\
\midrule
Endpoint validity & Confirmatory endpoint meets its pre-analysis definition and provenance rules. \\
Temporal transport & Pathway-aware model meets registered performance and calibration criteria in the locked later-period split. \\
Site transport & The result persists on an independent-site holdout or is explicitly downgraded as site-specific. \\
Robustness & Measurement-era, composition, missing-data, and block-ablation analyses do not negate the claim. \\
Genetic caution & Optional inherited block passes structure, relatedness, transferability, and subgroup-calibration review; otherwise it is removed. \\
Governance & Consent, privacy, specialist, statistical, and communications reviews approve the wording and release. \\
\bottomrule
\end{tabular}
\end{table}

\section{Expected Outcome Branches and Conditional Conclusions}

The study is planned so that several outcomes are informative. The first
branch is \emph{no incremental predictive value}. If the pathway-aware model
does not improve temporal and external-site calibration over baseline, the
result shows that the available records do not support an added pathway-aware
forecast. It does not show that pathways are biologically irrelevant. The
limitation may be poor measurement, inappropriate time alignment, limited
exposure range, or an endpoint that is not transportable.

The second branch is \emph{baseline explains apparent cohort trend}. Suppose
an unadjusted cohort association becomes weak or disappears after measurement
era and demographic-composition adjustment. Then the appropriate conclusion
is that the observed difference was not robust to non-biological competing
explanations. It must not be described as evidence against evolution in all
contexts; it is a result about this phenotype, dataset, and model.

The third branch is \emph{developmental or technology block improves
prediction}. If an exposure or treatment block adds stable, calibrated
predictive value, the conclusion may be that the block carries useful
population-level information in the eligible setting. The conclusion remains
associational. A separate causal design would be needed before claiming that
changing an exposure or treatment changes a phenotype.

The fourth branch is \emph{optional genetic block is unstable}. This is an
expected and useful possibility. If performance varies strongly with
population-structure controls, site, relatedness, or subgroup, the block is
not used in a forecast. That outcome implements, rather than merely cites,
the caution raised by Berg \emph{et al.} \cite{berg2019}.

The fifth branch is \emph{limited, transportable pathway attribution}. If
every gate passes for a clearly defined endpoint, the study may report a
conditional statement such as: in the eligible populations and measured
periods, a specified combination of measured pathway blocks improves
calibrated description of the phenotype over the baseline. It still cannot
predict an individual's appearance, assert inevitability across all humans, or
imply a hierarchy among human groups.

\section{Risks, Limitations, and Human Review Requirements}

\subsection{Scientific and statistical limitations}

The central scientific limitation is incomplete identification. Multiple
pathways can co-vary: an exposure can differ by site and period; treatment
availability can change with social conditions; and measured composition can
be an imperfect proxy for unmeasured history. The model comparison is useful
for prediction and competing-explanation audit, but it does not solve causal
identification. A causal effect of diet, technology, or an environmental
exposure would require a distinct target trial emulation, natural experiment,
or other justified design that is outside this proposal.

Measurement remains a decisive limitation. Craniofacial morphology is
particularly sensitive to landmark definitions, imaging method, posture, age,
and selection pathway. Dental records can be referral-biased and often lack
complete extraction history. An apparently large, multisite dataset can still
be less informative than a smaller, well-provenanced dataset. The protocol
therefore accepts a conservative result or no confirmatory analysis when
harmonization cannot be demonstrated.

Genetic interpretation is also constrained. An optional inherited-feature
summary may be less portable across cohorts, may encode residual population
structure, and may fail to improve calibration. It is not a measure of human
worth, biological identity, or a personal destiny. No individual-level result
is returned, and no model is deployed for clinical, educational, employment,
insurance, reproductive, policing, or immigration decisions.

\subsection{Ethics, privacy, and communication}

Human physical traits and genomic information are sensitive. The study
requires research-ethics approval, a lawful data-use basis, data minimization,
and a reidentification assessment. Facial images are excluded from released
analysis artifacts. Any linkage of phenotype and genotype uses approved,
pseudonymized records in a controlled environment. Access logging, retention
limits, and a process for withdrawal or custodial restrictions must be agreed
with the data holder before analysis.

Communication is not a cosmetic final step. A result that is technically well
calibrated can still be misrepresented as a story about racial destiny or a
universal future body. The release gate therefore includes review by domain
specialists, quantitative-methods reviewers, data-governance personnel, and
reviewers qualified to identify deterministic or stigmatizing language. The
public-facing summary must distinguish association, prediction, causal
inference, and inherited evolution.

\begin{table*}[t]
\caption{Human review responsibilities and stop conditions}
\label{tab:review}
\centering
\footnotesize
\begin{tabular}{p{0.20\textwidth}p{0.35\textwidth}p{0.35\textwidth}}
\toprule
Review role & Required decision & Stop or downgrade condition \\
\midrule
Research ethics and data governance & Confirm lawful basis, consent scope, minimization, retention, and reidentification controls. & Consent or data-use basis is unclear; raw images or identifiable genetic outputs are required. \\
Dental and craniofacial specialists & Approve endpoint manual, age adequacy, and measurement harmonization. & Agenesis cannot be separated from extraction or treatment; landmark definitions are not comparable. \\
Population-genetic specialist & Review genotype QC, structure, relatedness, transferability, and interpretation. & Optional inherited block remains structurally confounded or harms calibration in a supported stratum. \\
Statistician and reproducibility reviewer & Approve split logic, missing-data plan, metrics, code review, and preregistration. & Holdout leakage, unregistered target change, or unresolvable performance instability. \\
Communications and equity reviewer & Approve language and suppress disallowed inferences. & Wording invites individual prediction, population ranking, biological determinism, or eugenic interpretation. \\
\bottomrule
\end{tabular}
\end{table*}

\subsection{Scope boundary}

This proposal does not test whether humans will look different in a singular,
inevitable sense. It provides a way to ask narrower questions responsibly.
For each phenotype, the protocol can distinguish evidence for a measured
population-level association from a claim about inherited evolution, and it
can identify when the available evidence is too weak to support a forecast.
That is a strength, not an evasion: uncertainty is part of the scientific
result.

\section{Conclusion}

Future human physical phenotype cannot be responsibly inferred from a single
narrative about evolution, diet, climate, technology, or wisdom teeth.
Different generative pathways can yield similar visible patterns while
implying different scientific and ethical conclusions. The Phenotype Pathway
Attribution Atlas converts that recognition into a testable research design.
It defines bounded endpoints, separates congenital dental status from
treatment history, uses a cohort-and-measurement baseline, adds pathway blocks
conservatively, and requires temporal and external-site calibration before any
conditional forecast.

The proposal is intentionally strict about genetics and communication.
Inherited summaries are optional, consented, and sensitivity-tested; they are
never a population ranking or an individual appearance device. Technology and
developmental exposures are interpreted as associations unless a separate
causal design is justified. If the study fails its validation or governance
gates, it reports that failure without inventing a future-human story. If it
succeeds, its conclusion remains local to the measured phenotype, eligible
populations, periods, and modeling assumptions. This is the appropriate
scientific answer to a broad question: physical differences may occur, but
their pathways must be measured rather than imagined.

\appendices
\section{Idea-Evolution Record}

The workflow required a Survey--Idea--ExperimentDesign--Author provenance
chain. The Survey froze evidence and three gaps before idea selection.
IDEA-FP-01, a secular trend catalogue, was retained as a descriptive
alternative but rejected as the primary direction because it cannot separate
the relevant pathways. IDEA-FP-02, the Phenotype Pathway Attribution Atlas,
was selected because it ties each accepted gap to a measurable variable block
and a falsification condition. IDEA-FP-03, a universal future-human visual
generator, was rejected because it has no defensible universal target and
would create deterministic and stereotype risks.

No autonomous search trajectory, MCTS trace, experiment, simulation, or
multi-agent debate was executed. This appendix records the actual route
selection rather than fabricating an optimization history. The selected idea
remains a proposal: its central hypothesis is not that it has already improved
any metric, but that its benefit can be tested using locked temporal and
external-site holdouts.

\section{Symbols and Operational Definitions}

\begin{table}[h]
\caption{Symbols and definitions}
\label{tab:symbols}
\centering
\footnotesize
\begin{tabular}{p{0.22\columnwidth}p{0.68\columnwidth}}
\toprule
Symbol or term & Definition \\
\midrule
$Y_{ijk}$ & Phenotype $k$ for person $i$ at visit or age interval $j$. \\
$A_i$ & Confirmatory binary indicator of documented congenital third-molar agenesis, defined in \eqref{eq:agenesis}. \\
$\eta^{(0)}_{ijk}$ & Baseline model linear predictor containing age or stage, cohort, site, and measurement process. \\
$\eta^{(1)}_{ijk}$ & Pathway-aware model linear predictor in \eqref{eq:pathway}. \\
$\boldsymbol{x}_{ij}$ & Demographic-composition predictor block. \\
$\boldsymbol{z}_{ij}$ & Time-aligned developmental-exposure predictor block. \\
$\boldsymbol{q}_{ij}$ & Technology or treatment predictor block. \\
$\boldsymbol{r}_{i}$ & Optional separately consented inherited-feature summary. \\
Calibration & Agreement between predicted probabilities or intervals and observed frequencies or outcomes in held-out data. \\
Transportability & Preservation of registered predictive performance across later periods or independent sites. \\
\bottomrule
\end{tabular}
\end{table}

\section{Evidence Coverage and Review Checklist}

Table~\ref{tab:coverage} shows the allowed evidence role for each frozen
source. Literature sources motivate only the claims shown in the table; they
do not provide observed results for this study. The Author stage adds no
citations outside the Survey registry and no acknowledgements, personal
contributions, simulations, or result claims.

\begin{table*}[t]
\caption{Evidence coverage and Author-stage constraint check}
\label{tab:coverage}
\centering
\footnotesize
\begin{tabular}{p{0.20\textwidth}p{0.33\textwidth}p{0.37\textwidth}}
\toprule
Source & Frozen evidence role & Permitted Author-stage use \\
\midrule
Ponzi \emph{et al.} \cite{ponzi2020} & Developmental plasticity during ontogeny. & Separate developmental context from inherited-change claims. \\
Carter \emph{et al.} \cite{carter2015} & Heterogeneous prevalence of third-molar agenesis. & Justify a bounded dental endpoint and prohibit trend or selection overclaim. \\
Laland \emph{et al.} \cite{laland2010} & Gene--culture coevolution framework. & Model culture and technology as conditional environmental pathways. \\
Berg \emph{et al.} \cite{berg2019} & Population-stratification risk in polygenic adaptation inference. & Require conservative genetic sensitivity analysis. \\
1000 Genomes Project \cite{auton2015} & Human genetic variation resource. & Require QC and documented structure for any optional genetic block. \\
Collins \emph{et al.} \cite{collins2015} & Transparent prediction-model reporting. & Lock target, predictors, validation, and reporting. \\
Wolff \emph{et al.} \cite{wolff2019} & Prediction-model bias and applicability domains. & Audit participants, predictors, outcome, and analysis. \\
\bottomrule
\end{tabular}
\end{table*}

\begin{thebibliography}{00}

\bibitem{ponzi2020}
D. Ponzi, M. V. Flinn, M. P. Muehlenbein, and P. A. Nepomnaschy,
"Hormones and human developmental plasticity," \emph{Molecular and Cellular
Endocrinology}, vol. 505, Art. no. 110721, 2020, doi:
10.1016/j.mce.2020.110721.

\bibitem{carter2015}
K. Carter, C. Worthington, and J. M. A. M. Lewis, "Morphologic and
demographic predictors of third molar agenesis: A systematic review and
meta-analysis," \emph{Journal of Dental Research}, 2015, doi:
10.1177/0022034515581644.

\bibitem{laland2010}
K. N. Laland \emph{et al.}, "How culture shaped the human genome:
bringing genetics and the human sciences together," \emph{Nature Reviews
Genetics}, vol. 11, pp. 137--148, 2010, doi: 10.1038/nrg2734.

\bibitem{berg2019}
J. J. Berg \emph{et al.}, "Reduced signal for polygenic adaptation of
height in UK Biobank," \emph{eLife}, vol. 8, Art. no. e39725, 2019, doi:
10.7554/eLife.39725.

\bibitem{auton2015}
1000 Genomes Project Consortium, A. Auton \emph{et al.}, "A global
reference for human genetic variation," \emph{Nature}, vol. 526, pp.
68--74, 2015, doi: 10.1038/nature15393.

\bibitem{collins2015}
G. S. Collins \emph{et al.}, "Transparent reporting of a multivariable
prediction model for individual prognosis or diagnosis (TRIPOD): the TRIPOD
statement," \emph{Annals of Internal Medicine}, vol. 162, pp. 55--63, 2015,
doi: 10.7326/M14-0697.

\bibitem{wolff2019}
R. F. Wolff \emph{et al.}, "PROBAST: A tool to assess the risk of bias and
applicability of prediction model studies," \emph{Annals of Internal
Medicine}, vol. 170, pp. 51--58, 2019, doi: 10.7326/M18-1376.

\end{thebibliography}

\end{document}
